\section{Related Literature}\label{sec:literature}

The paper sits at the intersection of nonlinear service pricing, dynamic adoption and commitment, platform hold-up, and the emerging economics of AI and cloud services. The contribution is not the existence of any one of these ingredients, but the state-regime feedback created when sunk integration changes the state on which a later tariff-architecture decision is evaluated.

First, the nonlinear-pricing literature already gives a rich account of fixed fees, usage charges, and transaction costs. \citet{sundararajan2004} studies information goods when unlimited-usage and usage-based pricing are both feasible and administering usage pricing can be costly. That work is especially close to the static architecture choice in our continuation game. Our result does not challenge that analysis; instead, it asks what happens when the population exposed to the later pricing decision is itself selected by earlier expectations about the architecture.

Second, dynamic service-pricing models already show that commitment and forward-looking adoption can interact. \citet{penmetsa2015} studies dynamic pricing of new subscription services with strategic customers and explicitly emphasizes commitment power and technology adoption. \citet{gans2012} shows that when platform access is sunk before later application pricing, subsequent pricing can create an unravelling problem. These papers rule out a generic novelty claim of the form ``lack of commitment reduces adoption.'' In our model the relevant additional step is reciprocal: the adoption/integration response feeds back into the provider's later architecture incentive and splits one static switching condition into two self-consistency thresholds.

Third, tariff form has already been used as a commitment device in platform settings. \citet{mutherswismer2022} study sellers that incur platform-specific costs before a platform operator may enter their product markets; proportional fees can mitigate the resulting hold-up relative to classical two-part tariffs. This is a major conceptual neighbor. The opportunistic margin there is platform competition with sellers, whereas here it is the provider's later choice of a positive marginal usage price after observing the installed composition. We therefore do not claim novelty for tariff commitment or relationship-specific investment by themselves.

Fourth, recent AI and cloud papers substantially narrow any domain-based novelty claim. \citet{bergemannbonattismolin2025} model LLM pricing with variable token costs, fine-tuning, multidimensional user heterogeneity, and menus implementable through two-part tariffs. \citet{bergemannwang2025} study dynamic cloud contracts under uncertain future multi-unit demand and characterize implementations using two-part tariffs and committed-spend contracts. \citet{bhaskaran2026} analyze prepaid digital services with nontrivial serving costs, including AI and cloud examples, and examine extensions with two-part tariffs and endogenous quality. \citet{bichuchyaish2026} study a dynamic GenAI service in which free use is costly but can generate data that improves future quality. These papers establish that costly inference, tariff design, dynamic demand, and endogenous quality are already active research areas.

Against this background, the paper's claim is intentionally narrow. In the frozen baseline, anticipated metering lowers high-use integration while a larger high-use installed state raises the provider's gain from metering. This reciprocal ordering creates $\mu_M<\mu_F$ and therefore a region in which neither pure architecture is self-consistent. Reconstructing this result from the closest prior work would require combining a sunk integration decision with a later architecture choice whose payoff is itself altered by the integration state. The paper does not claim that this is the only possible dynamic pricing mechanism for generative AI, nor that it dominates richer screening or committed-spend mechanisms.
